\section{Montagens e densidade em espaços de funções contínuas}
\begin{definicao}[montagem contínua]
Se $K=[\alpha,\beta]$ é um intervalo compacto não degenerado, $\Asm(K)$ é o conjunto das funções $g\in C(K)$ para as quais existe uma coleção finita $\Lambda$ de células diádicas, com interiores disjuntos e união contendo $K$, e endereços $n(J)$, $J\in\Lambda$, tais que
\eqn{20}{g|_{J\cap K}=\ell_{n(J)}|_{J\cap K},}
com valores coincidentes nas fronteiras compartilhadas.
\end{definicao}
Para $f\in C(K)$, denote o módulo de continuidade por
\eqn{21}{\omega(f;\delta)=\sup\{|f(x)-f(y)|:x,y\in K,\ |x-y|\le\delta\}.}
O resultado seguinte impõe simultaneamente nós diádicos, valores racionais e inclinações não nulas. A figura~\ref{fig:montagem} mostra a ligação das peças por nós comuns.
\figura{03_montagem_continua}{1}{Montagem contínua de átomos. À esquerda, cada cor identifica uma célula e seu segmento; peças consecutivas compartilham o mesmo valor no nó. À direita, o refinamento substitui células por filhos e aproxima a função-alvo com mais peças conectadas.}{montagem}

\begin{teorema}[densidade das montagens]
Para todo $f\in C[\alpha,\beta]$ e todo $\varepsilon>0$, existe $g\in\Asm([\alpha,\beta])$ tal que
\[
\normK{f-g}{[\alpha,\beta]}<\varepsilon.
\]
Todas as peças de $g$ pertencem a $\LL$ e têm inclinação não nula.
\end{teorema}
\begin{proof}
Estenda $f$ por constantes: $F(x)=f(\alpha)$ para $x<\alpha$, $F=f$ em $[\alpha,\beta]$ e $F(x)=f(\beta)$ para $x>\beta$. A extensão é uniformemente contínua. Escolha $k$ de modo que, com $h=2^{-k}$,
\eqn{22}{\omega(F;h)<\varepsilon/2.}
Tome os nós $t_j=jh$ desde $j_{\min}=\lfloor\alpha/h\rfloor$ até $j_{\max}=\lceil\beta/h\rceil$. Para cada $j$, escolha racional $v_j$ com $|v_j-F(t_j)|<\varepsilon/2$ e, sequencialmente, imponha $v_{j+1}\ne v_j$. Isso é possível porque cada intervalo aberto de escolha contém infinitos racionais e apenas um valor é proibido em cada passo.

Seja $g$ o interpolante linear dos pares $(t_j,v_j)$. Em $J_{k,j}$, sua inclinação é $(v_{j+1}-v_j)/h\in\Q^*$ e seu valor no centro é $(v_j+v_{j+1})/2\in\Q$. Pelo teorema~3.6, a peça é o único átomo com esses parâmetros. Os valores nodais compartilhados tornam $g$ contínua.

Se $P_kF$ é o interpolante dos valores exatos, então, para $x=(1-\theta)t_j+\theta t_{j+1}$,
\[
|F(x)-P_kF(x)|\le(1-\theta)|F(x)-F(t_j)|+\theta|F(x)-F(t_{j+1})|
\le\omega(F;h)<\varepsilon/2.
\]
Além disso, $P_kF-g$ é linear em cada célula e seus valores nodais têm módulo menor que $\varepsilon/2$, logo $\norm{P_kF-g}<\varepsilon/2$. A desigualdade triangular conclui.
\end{proof}
O número de células dessa cobertura uniforme é
\eqn{23}{N_k=\lceil\beta2^k\rceil-\lfloor\alpha2^k\rfloor=(\beta-\alpha)2^k+O(1).}

\subsection{Equivalência exata com redes ReLU racionais}
Escreva $\rho(z)=z_+=\max\{0,z\}$. A identidade em hinges abaixo conserva os nós diádicos, os parâmetros racionais e as inclinações não nulas da montagem.
\begin{teorema}[equivalência montagem--ReLU racional]
Seja $g\in\Asm(K)$ dada por $N$ folhas consecutivas, com nós $t_0<t_1<\cdots<t_N$ e inclinação racional não nula $s_j$ em $[t_j,t_{j+1}]\cap K$, $0\le j<N$. Então, em $K$,
\eqn{24}{g(x)=g(t_0)+s_0(x-t_0)+\sum_{j=1}^{N-1}(s_j-s_{j-1})\rho(x-t_j).}
Depois de expandido o primeiro termo afim, todos os coeficientes são racionais, todos os pontos de quebra internos são diádicos e \eqref{eq:24} é uma rede ReLU rasa normalizada com $N-1$ unidades hinge e um termo afim de saída. Unidades cujo peso de saída é zero podem ser removidas.

Reciprocamente, seja
\eqn{25}{R(x)=c+qx+\sum_{j=1}^Ma_j\rho(x-\tau_j),\qquad c,q,a_j\in\Q,\quad\tau_j\in\D,}
com nós distintos. Suponha que $K$ admita uma cobertura por células diádicas compatível com os $\tau_j$, que os valores nodais sejam racionais e que cada inclinação local
\eqn{26}{q+\sum_{\tau_j<x}a_j}
seja não nula. Então $R|_K\in\Asm(K)$. Uma malha diádica uniforme que refine todos os $\tau_j$ fornece essa montagem explicitamente.
\end{teorema}
\begin{proof}
Na primeira folha, o lado direito de \eqref{eq:24} tem valor $g(t_0)$ em $t_0$ e inclinação $s_0$. Ao atravessar $t_j$, a derivada ganha exatamente $s_j-s_{j-1}$; por indução, sua inclinação na folha $j$ é $s_j$. Como ambos os lados são contínuos e coincidem no primeiro nó, coincidem em todo $K$. Cada $s_j$ é racional. Os nós internos são fronteiras de células diádicas, e o intercepto do primeiro átomo é racional por \eqref{eq:15}; logo a forma expandida tem somente parâmetros racionais.

Para a recíproca, refine uma grade diádica até que cada $\tau_j$ seja uma fronteira. Em qualquer célula aberta da grade, \eqref{eq:25} é afim, com inclinação \eqref{eq:26}. O valor no centro da célula é racional, pois o centro, os nós e todos os coeficientes são racionais. A hipótese de inclinação não nula permite aplicar o teorema~3.6, que fornece um único átomo da biblioteca nessa célula. A continuidade de $R$ faz os valores coincidirem nas fronteiras. A união das peças é, portanto, uma montagem no sentido da definição~4.1.
\end{proof}
A forma normalizada \eqref{eq:25} já implica valores nodais racionais. Se a arquitetura não oferecer conexão afim direta, $x=\rho(x)-\rho(-x)$ realiza o termo linear com duas unidades adicionais e viés constante na saída.

\begin{proposicao}[parâmetros e altura]
Na forma normalizada, a conversão de uma montagem com $N$ folhas armazena $2N$ racionais: o intercepto e a inclinação iniciais, mais um nó e um salto de inclinação para cada uma das $N-1$ unidades. Se todos os átomos têm altura no máximo $H\ge2$, a altura máxima dos parâmetros produzidos satisfaz
\eqn{27}{H_{\rm ReLU}\le2^{H+1}H^2.}
\end{proposicao}
\begin{proof}
Há $N-1$ pares $(t_j,s_j-s_{j-1})$, além dos dois parâmetros afins. Pela definição de altura atômica, $\heightq(s_j)\le H$, os nós satisfazem $\heightq(t_j)\le2^H$, e a forma geral do lema~7.11 dá $\heightq(s_j-s_{j-1})\le2H^2$. Para o intercepto inicial, escreva $a=y_0-sx_0$. Como $\heightq(x_0)\le2^H$,
\[
\heightq(a)\le2\heightq(y_0)\heightq(sx_0)\le2^{H+1}H^2.
\]
Isso prova \eqref{eq:27}. Como os demais parâmetros são produtos ou somas dos mesmos dados, essa cota também garante que a conversão é efetiva em aritmética racional exata.
\end{proof}
\figura{04_relu}{1}{Uma montagem racional com quatro folhas e a rede ReLU rasa dada pela identidade \eqref{eq:24}. O gráfico e os pesos são calculados a partir dos mesmos dados racionais.}{relu}
A identidade funcional do teorema~4.3 conduz aos formatos da seção~7.3, que distinguem a transmissão estrutural dos diádicos de sua expansão em frações reduzidas.

\subsection{Relação com o 1-jato de Taylor}
Se $J_{k,m}$ tem centro $x_0$ e largura $h$, o 1-jato de $f\in C^2(J_{k,m})$ é
\[
T_{x_0}f(x)=f(x_0)+f'(x_0)(x-x_0).
\]
\begin{proposicao}[pertencimento e aderência]
Se $f(x_0)\in\Q$ e $f'(x_0)\in\Q^*$, então $T_{x_0}f$ é o átomo de $\LL$ associado à célula. Em geral, para todo $\delta>0$ existem $q_0\in\Q$ e $q_1\in\Q^*$ tais que $\ell_q(x)=q_0+q_1(x-x_0)$ satisfaz
\eqn{28}{\normK{T_{x_0}f-\ell_q}{J_{k,m}}<\delta,}
\eqn{29}{\normK{f-\ell_q}{J_{k,m}}\le\frac{h^2}8\normK{f''}{J_{k,m}}+\delta.}
\end{proposicao}
\begin{proof}
No primeiro caso, $(x_0,f(x_0),f'(x_0))$ pertence ao contradomínio do teorema~3.6. No caso geral, escolha $q_0$ com $|f(x_0)-q_0|<\delta/2$ e $q_1\ne0$ com $|f'(x_0)-q_1|<\delta/h$. Como $|x-x_0|\le h/2$, obtém-se \eqref{eq:28}. O resto integral de Taylor fornece
\[
|f(x)-T_{x_0}f(x)|=\left|\int_{x_0}^x(x-t)f''(t)\,dt\right|
\le\frac{|x-x_0|^2}2\norm{f''}\le\frac{h^2}8\norm{f''}.
\]
A desigualdade triangular prova \eqref{eq:29}.
\end{proof}
Jatos reais arbitrários pertencem, portanto, à aderência local da biblioteca racional.
